\documentclass{stex}

\libusepackage{naproche}
\libinput{preamble}

\begin{document}
\begin{smodule}{euclid-division-theorem.ftl}

\importmodule[libraries/arithmetics]{multiplication-and-ordering.ftl}
\importmodule[libraries/arithmetics]{subtraction.ftl}

\symdecl*{Euclid Division Existence}

\begin{theorem}[forthel,title=Euclid's Division Theorem: Existence,name=Euclid Division Existence]
  Let $n, m$ be natural numbers such that $m \NotEq \Zero$.
  Then there exist natural numbers $q, r$ such that
  \[n \Eq (m \Mul q) \Plus r\]
  and $r \Less m$.
\end{theorem}
\begin{proof}[forthel]
  (a) Define $\Phi \DefEq \SClass{n'}{\Naturals}{\text{ there exist natural numbers }q, r\text{ such that }r \Less m\text{ and }n' \Eq (m \Mul q) \Plus r}$.

  (1) $\Phi$ contains $\Zero$.
  \begin{proof}
    Take $q \Eq \Zero$ and $r \Eq \Zero$.
    Then $r \Less m$ and $\Zero \Eq (m \Mul q) \Plus r$.
    Hence $\Zero \Elem \Phi$.
  \end{proof}

  (2) For all $n' \Elem \Phi$ we have $n' \Plus \One \Elem \Phi$.
  \begin{proof}
    Let $n' \Elem \Phi$.

    Let us show that there exist natural numbers $q, r$ such that $r \Less m$ and $n' \Plus \One \Eq (m \Mul q) \Plus r$.
      Take natural numbers $q', r'$ such that $r' \Less m$ and $n' \Eq (m \Mul q') \Plus r'$ (by a).
      We have $r' \Plus \One \Less m$ or $r' \Plus \One \Eq m$.

      \begin{case}{$r' \Plus \One \Less m$.}
        Take $q \Eq q' \Plus \Zero$ and $r \Eq r' \Plus \One$. %!
        Then $r \Less m$ and $n' \Plus \One
          \Eq ((q' \Mul m) \Plus r') \Plus \One
          \Eq (q' \Mul m) \Plus (r' \Plus \One)
          \Eq (q \Mul m) \Plus r$.
      \end{case}

      \begin{case}{$r' \Plus \One \Eq m$.}
        Take  $q \Eq q' \Plus \One$ and $r \Eq \Zero$.
        Then $r \Less m$ and
        $n' \Plus \One
          \Eq ((q' \Mul m) \Plus r') \Plus \One
          \Eq (q' \Mul m) \Plus (r' \Plus \One)
          \Eq (q' \Mul m) \Plus m
          \Eq (q' \Mul m) \Plus (\One \Mul m)
          \Eq (q' \Plus \One) \Mul m
          \Eq (q \Mul m) \Plus r$.
      \end{case}
    End.

    Hence $n' \Plus \One \Elem \Phi$.
  \end{proof}

  Then $\Phi$ contains every natural number (by \sr{induction I}{induction}).
  Thus there exist natural numbers $q, r$ such that $n \Eq (m \Mul q) \Plus r$ and $r \Less m$ (by a).
\end{proof}

\symdecl*{Euclid Division Uniqueness}

\begin{theorem}[forthel,title=Euclid's Division Theorem: Uniqueness,name=Euclid Division Uniqueness]
  Let $n, m$ be natural numbers such that $m \NotEq \Zero$.
  Let $q, r$ be natural numbers such that
  \[n \Eq (m \Mul q) \Plus r\]
  and $r \Less m$.
  Let $q', r'$ be natural numbers such that
  \[n \Eq (m \Mul q') \Plus r'\]
  and $r' \Less m$.
  Then $q \Eq q'$ and $r \Eq r'$.
\end{theorem}
\begin{proof}[forthel]
  We have $(m \Mul q) \Plus r \Eq (m \Mul q') \Plus r'$.

  \begin{case}{$q \Geq q'$ and $r \Geq r'$.}
    (1) $((m \Mul q) \Plus r) \Minus r' \Eq (m \Mul q) \Plus (r \Minus r')$ (by \sn{associativity of addition and subtraction}).
    (2) $((m \Mul q') \Plus r') \Minus r'
      \Eq (m \Mul q') \Plus (r' \Minus r')
      \Eq m \Mul q'$.
    Hence $(m \Mul q) \Plus (r \Minus r') \Eq m \Mul q'$.
    Thus $((m \Mul q) \Minus (m \Mul q')) \Plus (r \Minus r') \Eq \Zero$.
    Consequently $(m \Mul q) \Minus (m \Mul q') \Eq \Zero$ and $r \Minus r' \Eq \Zero$.
    If $(m \Mul q) \Minus (m \Mul q') \Eq \Zero$ then $q \Minus q' \Eq \Zero$.
    Therefore $q \Minus q' \Eq \Zero$ and $r \Minus r' \Eq \Zero$.
    Thus we have $q \Eq q'$ and $r \Eq r'$.
  \end{case}

  \begin{case}{$q \Geq q'$ and $r \Less r'$.}
    Take $q'' \Eq q \Minus q'$ and $r'' \Eq r' \Minus r$.
    Then $(m \Mul (q' \Plus q'')) \Plus r \Eq (m \Mul q') \Plus (r \Plus r'')$.
    We have $(m \Mul q') \Plus (r \Plus r'')
      \Eq (m \Mul q') \Plus (r'' \Plus r)
      \Eq ((m \Mul q') \Plus r'') \Plus r$.
    Hence $(m \Mul (q' \Plus q'')) \Plus r \Eq ((m \Mul q') \Plus r'') \Plus r$.
    Thus $m \Mul (q' \Plus q'') \Eq (m \Mul q') \Plus r''$ (by \sn{right-cancellability of addition}).
    We have $m \Mul (q' \Plus q'') \Eq (m \Mul q') \Plus (m \Mul q'')$.
    Hence $(m \Mul q') \Plus (m \Mul q'') \Eq (m \Mul q') \Plus r''$.
    [prover vampire]
    Thus $m \Mul q'' \Eq r''$.
    Then we have $m \Mul q'' \Less m \Mul \One$.
    Indeed $m \Mul q''
      \Eq r''
      \Leq r'
      \Less m
      \Eq m \Mul \One$.
    Therefore $q'' \Less \One$ (by \sn{preservation of ordering under left-multiplication}).
    Consequently $q \Minus q' \Eq q'' \Eq \Zero$.
    Hence $q \Eq q'$.
    Thus $(m \Mul q) \Plus r \Eq (m \Mul q) \Plus r'$.
    Therefore $r \Eq r'$.
  \end{case}

  \begin{case}{$q \Less q'$ and $r \Geq r'$.}
    Take $q'' \Eq q' \Minus q$ and $r'' \Eq r \Minus r'$.
    Then $(m \Mul q) \Plus (r' \Plus r'') \Eq (m \Mul (q \Plus q'')) \Plus r'$.
    We have $(m \Mul q) \Plus (r' \Plus r'')
      \Eq (m \Mul q) \Plus (r'' \Plus r')
      \Eq ((m \Mul q) \Plus r'') \Plus r'$.
    Hence $((m \Mul q) \Plus r'') \Plus r' \Eq (m \Mul (q \Plus q'')) \Plus r'$.
    Thus $(m \Mul q) \Plus r'' \Eq m \Mul (q \Plus q'')$ (by \sn{right-cancellability of addition}).
    We have $m \Mul (q \Plus q'') \Eq (m \Mul q) \Plus (m \Mul q'')$.
    Hence $(m \Mul q) \Plus r'' \Eq (m \Mul q) \Plus (m \Mul q'')$.
    [prover vampire]
    Thus $r'' \Eq m \Mul q''$.
    Then we have $m \Mul q'' \Less m \Mul \One$.
    Indeed $m \Mul q''
      \Eq r''
      \Leq r
      \Less m
      \Eq m \Mul \One$.
    Therefore $q'' \Less \One$ (by \sn{preservation of ordering under left-multiplication}).
    Consequently $q' \Minus q \Eq q'' \Eq \Zero$.
    Hence $q' \Eq q$.
    Thus $(m \Mul q) \Plus r \Eq (m \Mul q) \Plus r'$.
    Therefore $r \Eq r'$.
  \end{case}

  \begin{case}{$q \Less q'$ and $r \Less r'$.}
    (1) $((m \Mul q') \Plus r') \Minus r \Eq (m \Mul q') \Plus (r' \Minus r)$ (by \sn{associativity of addition and subtraction}).
    (2) $((m \Mul q) \Plus r) \Minus r
      \Eq (m \Mul q) \Plus (r \Minus r)
      \Eq m \Mul q$.
    Hence $(m \Mul q') \Plus (r' \Minus r) \Eq m \Mul q$.
    Thus $((m \Mul q') \Minus (m \Mul q)) \Plus (r' \Minus r)
      \Eq ((m \Mul q') \Plus (r' \Minus r)) \Minus (m \Mul q)
      \Eq \Zero$.
    Consequently $(m \Mul q') \Minus (m \Mul q) \Eq \Zero$ and $r' \Minus r \Eq \Zero$.
    If $(m \Mul q') \Minus (m \Mul q) \Eq \Zero$ then $q' \Minus q \Eq \Zero$.
    Therefore $q' \Minus q \Eq \Zero$ and $r' \Minus r \Eq \Zero$.
    Thus we have $q' \Eq q$ and $r' \Eq r$.
  \end{case}
\end{proof}
\end{smodule}
\end{document}
